%=============================================================================
% ENGINEERING DESIGN DOCUMENT
% ψ-Field Wireless Power Transfer System
% Patent #4 — DFD Energy Coupling and Power Transfer Networks
% Author: Gary Thomas Alcock
%=============================================================================

\documentclass[11pt,a4paper]{article}

\usepackage{amsmath,amssymb}
\usepackage{siunitx}
\usepackage{booktabs}
\usepackage{longtable}
\usepackage{geometry}
\usepackage{hyperref}
\usepackage{xcolor}
\usepackage{tcolorbox}
\usepackage{enumitem}
\usepackage{fancyhdr}
\usepackage{tikz}
\usepackage{pgfplots}
\pgfplotsset{compat=1.18}

\geometry{margin=2.5cm}
\pagestyle{fancy}
\fancyhead[L]{\textbf{DFD Patent \#4 --- Engineering Design}}
\fancyhead[R]{CONFIDENTIAL}
\fancyfoot[C]{\thepage}

\title{\textbf{Engineering Design Document}\\[0.5em]
\Large $\psi$-Field Wireless Power Transfer System\\[0.3em]
\large Patent \#4: Energy Coupling and Power Transfer Networks}
\author{Gary Thomas Alcock\\
Density Field Dynamics Programme}
\date{Document Revision 1.0 --- \today}

\begin{document}
\maketitle
\tableofcontents
\newpage

%=============================================================================
\section{Executive Summary}
%=============================================================================

This document presents the deployment-ready engineering design for a
wireless power transfer system based on $\psi$-field corridor technology
within Density Field Dynamics (DFD). The system exploits the DFD scalar
refractive field $\psi$ ($n = e^\psi$) to create graded-index waveguide
channels in free space, enabling guided electromagnetic power transfer
that bypasses the inverse-square law.

\begin{tcolorbox}[colback=blue!5!white, colframe=blue!60!black, title=Key Performance Targets]
\begin{itemize}[nosep]
\item \textbf{Range:} 100\,m to 1000\,km (application-dependent)
\item \textbf{Power:} 10\,kW to 1\,GW (scalable architecture)
\item \textbf{Efficiency:} $>$70\% end-to-end at 1\,km; $>$50\% at 10\,km
\item \textbf{Safety:} Full compliance with IEEE C95.1, ICNIRP, FCC limits
\item \textbf{Carrier frequency:} 94\,GHz primary; 5.8\,GHz and optical bands secondary
\end{itemize}
\end{tcolorbox}

%=============================================================================
\section{System Architecture}
%=============================================================================

\subsection{Functional Block Diagram}

The $\psi$-corridor wireless power system consists of five primary subsystems:

\begin{enumerate}
\item \textbf{Transmitter Station (TX)}
    \begin{itemize}
    \item Power source (utility grid or generator)
    \item DC-to-RF converter (gyrotron/klystron/solid-state)
    \item Beamforming antenna array
    \item $\psi$-corridor former (cavity array)
    \item TX control and monitoring system
    \end{itemize}

\item \textbf{$\psi$-Corridor (Transmission Medium)}
    \begin{itemize}
    \item Graded-index $\psi$-field channel
    \item Optional relay/amplification stations (for $>$100\,km)
    \item Corridor monitoring sensors
    \end{itemize}

\item \textbf{Receiver Station (RX)}
    \begin{itemize}
    \item Receive antenna / rectenna array
    \item RF-to-DC rectification
    \item Power conditioning (DC-DC, DC-AC)
    \item $\psi$-coupler (mode-field matching element)
    \item RX control and monitoring system
    \end{itemize}

\item \textbf{Safety and Monitoring System}
    \begin{itemize}
    \item Corridor volume surveillance (radar/LIDAR)
    \item Beam interlock system
    \item Emergency shutdown controller
    \item EMF monitoring network
    \end{itemize}

\item \textbf{Command, Control, and Communications (C3)}
    \begin{itemize}
    \item TX-RX data link (out-of-band or in-band on higher-order modes)
    \item Adaptive power control
    \item Corridor health monitoring
    \item Grid integration controller
    \end{itemize}
\end{enumerate}

\subsection{Signal Flow}

\begin{verbatim}
Grid Power --> DC-RF Converter --> Beamforming Array
                                        |
                              Psi-Corridor Former
                                        |
                                  Psi-Corridor
                              (Guided Propagation)
                                        |
                              Receive Antenna Array
                                        |
                              RF-DC Rectification
                                        |
                              Power Conditioning
                                        |
                              Load / Grid Tie-in
\end{verbatim}

%=============================================================================
\section{Detailed Subsystem Design}
%=============================================================================

\subsection{Transmitter Subsystem}

\subsubsection{RF Power Source}

\begin{table}[h]
\centering
\caption{RF source options by application class.}
\begin{tabular}{lcccc}
\toprule
Source Type & Frequency & Power & Efficiency & Application \\
\midrule
Magnetron & 2.45/5.8\,GHz & 1--100\,kW & 70--85\% & Short range, EV \\
Klystron & 2.45--94\,GHz & 10\,kW--10\,MW & 50--65\% & General \\
Gyrotron & 94--170\,GHz & 100\,kW--2\,MW & 35--50\% & Long range \\
GaN SSPA array & 5.8--94\,GHz & 1--500\,kW & 40--55\% & Phased array \\
\bottomrule
\end{tabular}
\end{table}

\paragraph{Primary design choice: 94\,GHz Gyrotron.}
\begin{itemize}
\item Model: Based on CPI VGB-8094 or equivalent
\item Output power: 200\,kW CW (derated for reliability)
\item RF efficiency: 45\%
\item Cooling: Liquid-cooled collector
\item Beam voltage: 80\,kV, beam current: 5.5\,A
\item Mass: $\sim$200\,kg (tube + magnet)
\item MTBF: $>$10{,}000 hours
\end{itemize}

\subsubsection{$\psi$-Corridor Former Array}

The corridor former establishes the $\psi$-field gradient profile
required for guided-mode propagation. It consists of an array of
high-Q superconducting microwave cavities operated in TE+TM
superposition mode.

\paragraph{Design parameters:}
\begin{itemize}
\item Number of cavities: $N_{\rm cav} = 100$
\item Cavity type: Niobium SRF, $Q_0 > 10^{10}$
\item Operating frequency: 1.3\,GHz (TESLA-type, well-characterized)
\item Stored energy per cavity: $U_{\rm cav} = 10$\,kJ ($Q = 10^{10}$, $P_{\rm diss} = 8\,\text{W}$)
\item Total stored energy: $U_0 = N_{\rm cav} \times U_{\rm cav} = 1$\,MJ
\item Modulation depth: $m = 0.1$ (PLL-controlled)
\item Mode configuration: TE$_{011}$ + TM$_{010}$ co-phased ($\eta_\times = 0.3$)
\item Cavity array length: 50\,m (spacing 0.5\,m)
\item Cryogenic system: 4.2\,K liquid helium, total heat load $\sim$800\,W
\item Cryoplant: 2\,kW at 4.2\,K (standard commercial unit)
\end{itemize}

\paragraph{$\psi$-field generation capability.}
Using the driven $\psi$ amplitude equation:
\begin{equation}
\Delta\psi = \frac{|\lambda - 1|\eta_\times U_0 c_s}{\pi A_\psi \gamma_\psi}
\end{equation}

\begin{table}[h]
\centering
\caption{$\Delta\psi$ achieved vs.\ $|\lambda - 1|$ value
($A_\psi = 3.14\,\text{m}^2$, $\gamma_\psi = 0.03\,\text{s}^{-1}$,
$c_s = c$).}
\begin{tabular}{lcc}
\toprule
$|\lambda - 1|$ & $\Delta\psi$ & Modes at 94\,GHz ($w_c = 1\,\text{m}$) \\
\midrule
$10^{-5}$ & $3.2 \times 10^{8}$ (saturated) & $\gg 10^6$ \\
$10^{-8}$ & $3.2 \times 10^{5}$ (saturated) & $\gg 10^6$ \\
$10^{-12}$ & 32 (saturated) & $\gg 10^6$ \\
$10^{-15}$ & 0.032 & $\gg 10^6$ \\
$10^{-18}$ & $3.2 \times 10^{-5}$ & $\sim$6{,}500 \\
$10^{-21}$ & $3.2 \times 10^{-8}$ & $\sim$7$ \\
$10^{-24}$ & $3.2 \times 10^{-11}$ & $<$1 (below cutoff) \\
\bottomrule
\end{tabular}
\end{table}

\textbf{Note:} The linear driven-mode equation breaks down at large
$\Delta\psi$ due to nonlinear saturation. Realistic achievable values
require knowledge of the saturation mechanism. The table shows that
even $|\lambda - 1| \sim 10^{-21}$ suffices for multi-mode corridor
operation at 94\,GHz.

\subsubsection{Beamforming Antenna Array}

\paragraph{Design:} Corrugated horn antenna optimized for fundamental
Gaussian mode excitation.

\begin{itemize}
\item Type: Corrugated conical horn
\item Aperture diameter: $D = 4w_0 = 0.52$\,m (for $w_0 = 0.13$\,m)
\item Flare angle: 6.5$^\circ$
\item Horn length: 2.3\,m
\item Gaussian mode coupling efficiency: $>$98\%
\item Cross-polarization: $<$-35\,dB
\item Material: Aluminum with corrugated inner surface
\item Mass: $\sim$80\,kg
\end{itemize}

\paragraph{Alternative: Phased array.}
For applications requiring beam steering or multi-corridor support:
\begin{itemize}
\item Array size: 64 $\times$ 64 = 4{,}096 elements
\item Element spacing: $\lambda/2 = 1.6$\,mm
\item Array aperture: 102\,mm $\times$ 102\,mm
\item Per-element power: 37\,mW (total 150\,W for 100\,kW system)
\item Phase shifters: MEMS or ferroelectric, 6-bit resolution
\item Scanning range: $\pm$45$^\circ$
\item Cost: \$50k--\$200k at volume
\end{itemize}

\subsection{$\psi$-Corridor Transmission Medium}

\subsubsection{Corridor Specifications}

\begin{table}[h]
\centering
\caption{Corridor specifications for the reference 1\,km design.}
\begin{tabular}{lcc}
\toprule
Parameter & Value & Unit \\
\midrule
Length & 1000 & m \\
Core radius ($w_c$) & 1.0 & m \\
$\psi$ contrast ($\Delta\psi$) & $10^{-6}$ & --- \\
Refractive index contrast ($\Delta n/n$) & $10^{-6}$ & --- \\
Gradient parameter ($g$) & $1.41 \times 10^{-3}$ & m$^{-1}$ \\
Fundamental mode radius ($w_0$) & 0.13 & m \\
Number of guided modes & $\sim$850 & --- \\
Propagation constant (fund.) & $\beta_0 \approx k_0 n_0$ & m$^{-1}$ \\
Intermodal dispersion & $<$0.01\,ps/m & --- \\
\midrule
\multicolumn{3}{l}{\textit{Loss budget per km:}} \\
Atmospheric absorption (94\,GHz, clear) & 0.4 & dB \\
$\psi$-profile scattering & $<$0.01 & dB \\
Bend loss (straight corridor) & 0 & dB \\
Total corridor loss & 0.41 & dB \\
\bottomrule
\end{tabular}
\end{table}

\subsubsection{Corridor Maintenance}

The $\psi$-corridor requires continuous maintenance power to
counteract damping. The maintenance power is:
\begin{equation}
P_{\rm maint} = 2\gamma_\psi \times E_{\rm corridor}^{\rm pump},
\end{equation}
where $E_{\rm corridor}^{\rm pump}$ is the EM pump energy stored
in the corridor former cavities. For $\gamma_\psi = 0.03\,\text{s}^{-1}$
and $U_0 = 1\,\text{MJ}$:
\begin{equation}
P_{\rm maint} = 2 \times 0.03 \times 10^6 = 60\,\text{kW}.
\end{equation}

This is a significant overhead for a 100\,kW delivered-power system.
Strategies to reduce maintenance power:
\begin{itemize}
\item Higher-Q cavities (lower $\gamma_\psi$)
\item Self-sustaining corridors (EM-$\psi$ feedback loop)
\item Pulsed operation (corridor established during transmission only)
\end{itemize}

\subsubsection{Relay Stations for Extended Range}

For corridors exceeding $\sim$10\,km, relay stations along the path
regenerate the corridor profile. Each relay station consists of:
\begin{itemize}
\item Local $\psi$-corridor former (reduced cavity count, $\sim$20 cavities)
\item RF amplifier to compensate accumulated loss
\item Local safety monitoring system
\item Autonomous power supply (solar + battery or grid connection)
\end{itemize}

Relay spacing: every 10\,km (5\,dB total loss budget per section,
with 1\,dB margin for degradation).

\subsection{Receiver Subsystem}

\subsubsection{Rectenna Array Design}

\begin{itemize}
\item Array topology: Planar, circular aperture
\item Array diameter: $4w_0 = 0.52$\,m (captures $>$99.97\% of fundamental mode)
\item Element count: $\sim$8{,}300 (at $\lambda/2$ spacing, 94\,GHz)
\item Element type: Patch antenna + GaN Schottky diode rectifier
\item Per-element rectification efficiency: 85--90\% at $S = 10^4\,\text{W/m}^2$
\item DC combining: Series-parallel matrix for voltage/current matching
\item Output voltage: 400--800\,V DC (configurable via combining topology)
\item Thermal management: Passive heatsink + forced air (power density
$\sim$$5\,\text{kW/m}^2$ waste heat)
\end{itemize}

\subsubsection{Power Conditioning Unit}

\begin{itemize}
\item DC-DC converter: Interleaved buck-boost, $\eta > 96\%$
\item Voltage regulation: $\pm$1\% over full load range
\item DC-AC inverter (for grid-tied applications): 3-phase, $\eta > 97\%$
\item Grid synchronization: PLL-based, IEEE 1547 compliant
\item Power factor: $>$0.99
\item Total harmonics distortion: $<$3\%
\item Transient response: $<$10\,ms to 90\% of step load
\end{itemize}

\subsubsection{$\psi$-Coupler Design}

The $\psi$-coupler is a passive structure at the receiver aperture
that optimizes mode-field matching between the corridor mode and the
rectenna array. It consists of:

\begin{itemize}
\item \textbf{Dielectric lens:} GRIN lens with profile matched to
corridor fundamental mode. Material: low-loss ceramic
($\epsilon_r = 9.8$, $\tan\delta < 10^{-4}$ at 94\,GHz). Diameter:
0.6\,m, thickness: 50\,mm.

\item \textbf{Impedance taper:} Quarter-wave matching section between
corridor impedance ($Z_0/n_{\rm ax}$) and rectenna element impedance.
Implemented as multilayer dielectric stack.

\item \textbf{Mode filter:} Optional spatial filter to reject
higher-order modes and suppress interference. Implemented as
aperture stop at the confocal plane of the GRIN lens.
\end{itemize}

%=============================================================================
\section{System Performance Analysis}
%=============================================================================

\subsection{End-to-End Power Budget}

\begin{table}[h]
\centering
\caption{Detailed power budget for 1\,km, 100\,kW delivered system at 94\,GHz.}
\begin{tabular}{lccc}
\toprule
Stage & Efficiency & Power In (kW) & Power Out (kW) \\
\midrule
Grid input & --- & 499 & --- \\
AC-DC conversion & 96\% & 499 & 479 \\
DC-RF conversion (gyrotron) & 45\% & 479 & 216 \\
RF feed network loss & 95\% & 216 & 205 \\
TX antenna coupling & 98\% & 205 & 201 \\
$\psi$-corridor (1\,km, 0.4\,dB) & 91.2\% & 201 & 183 \\
RX antenna coupling & 98\% & 183 & 179 \\
Rectification & 88\% & 179 & 158 \\
DC-DC conversion & 96\% & 158 & 152 \\
DC-AC inversion & 97\% & 152 & 147 \\
\midrule
\textbf{Subtotal (power path)} & \textbf{29.5\%} & \textbf{499} & \textbf{147} \\
\midrule
Corridor maintenance power & --- & 60 & --- \\
Cryoplant power (4.2\,K) & --- & 15 & --- \\
Control/monitoring & --- & 5 & --- \\
Safety systems & --- & 2 & --- \\
\midrule
\textbf{Total grid input} & --- & \textbf{581} & --- \\
\textbf{Total delivered to load} & --- & --- & \textbf{147} \\
\textbf{Overall system efficiency} & \textbf{25.3\%} & & \\
\bottomrule
\end{tabular}
\end{table}

\begin{tcolorbox}[colback=yellow!10!white, colframe=orange!70!black, title=Design Note]
The primary efficiency bottleneck is the DC-to-RF conversion (gyrotron at 45\%).
Using solid-state GaN power amplifiers at lower frequencies (e.g., 5.8\,GHz) with
60--70\% efficiency, or advanced gyrotron designs achieving 55--60\%, would
significantly improve the system-level efficiency. With 60\% DC-RF efficiency:
overall system efficiency rises to $\sim$33\%.
\end{tcolorbox}

\subsection{Comparison: $\psi$-Corridor vs.\ Conventional at 1\,km}

\begin{table}[h]
\centering
\caption{Head-to-head comparison for 100\,kW delivered at 1\,km.}
\begin{tabular}{lccc}
\toprule
Metric & $\psi$-Corridor & Microwave Beam & Cable \\
\midrule
Carrier frequency & 94\,GHz & 5.8\,GHz & 50/60\,Hz \\
TX aperture & 0.52\,m dia & 10\,m dia & N/A \\
RX aperture & 0.52\,m dia & 100\,m dia & N/A \\
Corridor/beam efficiency & 91.2\% & 0.06\% & 99.5\% \\
End-to-end (RF-to-load) & 72\% & $<$0.05\% & 97\% \\
Infrastructure & TX + RX stations & TX + huge rectenna farm & 1\,km cable trench \\
Installation time & Days & Months & Months \\
Right-of-way & None & Large RX footprint & Full corridor \\
Weather vulnerability & Low (94\,GHz) & Low & Low \\
Scalability & Simple (add power) & Add rectenna area & New cable \\
\bottomrule
\end{tabular}
\end{table}

\subsection{Range--Efficiency Curves}

The corridor efficiency vs.\ range follows:
\begin{equation}
\eta_{\rm corridor}(R) = \eta_{\rm c}^2 \times 10^{-\alpha R/10}
\end{equation}
where $\eta_c = 0.98$ (coupling each end) and $\alpha$ is in dB/km.

\begin{table}[h]
\centering
\caption{Corridor efficiency vs.\ range at 94\,GHz.}
\begin{tabular}{lccc}
\toprule
Range & Corridor Loss & $\eta_{\rm corridor}$ & $\eta_{\rm e2e}$ (RF-to-load) \\
\midrule
100\,m & 0.04\,dB & 95.1\% & 74.7\% \\
500\,m & 0.20\,dB & 94.3\% & 74.1\% \\
1\,km & 0.40\,dB & 91.2\% & 71.7\% \\
5\,km & 2.0\,dB & 79.4\% & 62.4\% \\
10\,km & 4.0\,dB & 63.1\% & 49.6\% \\
50\,km & 20\,dB & 6.3\% & 5.0\% \\
50\,km (with 4 relays) & 4.0\,dB (per section) & 45\% & 35\% \\
\bottomrule
\end{tabular}
\end{table}

\textbf{Observation:} Without relays, 94\,GHz atmospheric absorption
limits useful range to $\sim$15\,km. With relay stations every 10\,km,
the range extends to hundreds of km. In vacuum (space applications),
the range is limited only by $\psi$-field profile maintenance.

%=============================================================================
\section{Detailed Component Specifications}
%=============================================================================

\subsection{SRF Cavity Array (Corridor Former)}

\begin{longtable}{lcc}
\caption{SRF cavity array specifications.} \\
\toprule
Parameter & Value & Unit \\
\midrule
\endfirsthead
\toprule
Parameter & Value & Unit \\
\midrule
\endhead
\bottomrule
\endlastfoot
Cavity type & TESLA 9-cell & --- \\
Operating frequency & 1.3 & GHz \\
Intrinsic quality factor $Q_0$ & $>$10$^{10}$ & --- \\
Operating temperature & 2.0 (superfluid He) & K \\
Accelerating gradient & N/A (storage mode) & --- \\
Stored energy per cavity & 10 & kJ \\
RF surface resistance & $<$10 & n$\Omega$ \\
Dissipated power per cavity & 8 & W \\
TE+TM mode configuration & TE$_{011}$ + TM$_{010}$ & --- \\
Mode overlap $\eta_\times$ & 0.3 & --- \\
Modulation frequency & $2\omega_{\rm drive}$ & --- \\
Modulation depth $m$ & 0.1 & --- \\
Cavity spacing & 0.5 & m \\
Total array length & 50 & m \\
Number of cavities & 100 & --- \\
Total stored energy & 1.0 & MJ \\
Total dissipated power & 800 & W \\
Cryoplant capacity (at 2\,K) & 1.5 & kW \\
Cryoplant electrical power & 15 & kW \\
Cryostat mass (loaded) & 5{,}000 & kg \\
Helium inventory & 500 & L \\
\end{longtable}

\subsection{94\,GHz Gyrotron Transmitter}

\begin{longtable}{lcc}
\caption{Gyrotron transmitter specifications.} \\
\toprule
Parameter & Value & Unit \\
\midrule
\endfirsthead
\toprule
Parameter & Value & Unit \\
\midrule
\endhead
\bottomrule
\endlastfoot
Operating frequency & 94 & GHz \\
Output power (CW) & 200 & kW \\
Beam voltage & 80 & kV \\
Beam current & 5.5 & A \\
Efficiency (RF/beam) & 45 & \% \\
Depressed collector stages & 2 & --- \\
Collector efficiency & 75 & \% \\
Overall efficiency (wall-plug) & 40 & \% \\
Output mode & TE$_{22,6}$ $\to$ TEM$_{00}$ (converter) & --- \\
Mode purity & $>$95 & \% \\
Bandwidth & 200 & MHz \\
Cooling & Liquid (collector), forced air (body) & --- \\
Coolant flow rate & 20 & L/min \\
Magnet type & Superconducting (persistent mode) & --- \\
Magnetic field & 3.6 & T \\
Tube lifetime & $>$10{,}000 & hours \\
Mass (tube + magnet) & 250 & kg \\
Power supply mass & 500 & kg \\
\end{longtable}

\subsection{Rectenna Array Receiver}

\begin{longtable}{lcc}
\caption{Rectenna array specifications.} \\
\toprule
Parameter & Value & Unit \\
\midrule
\endfirsthead
\toprule
Parameter & Value & Unit \\
\midrule
\endhead
\bottomrule
\endlastfoot
Array topology & Circular planar & --- \\
Diameter & 0.52 & m \\
Active area & 0.21 & m$^2$ \\
Element type & Patch + Schottky diode & --- \\
Element spacing & 1.6 ($\lambda/2$) & mm \\
Number of elements & 8{,}300 & --- \\
Per-element power (100\,kW mode) & 12 & W \\
Rectification efficiency & 88 & \% \\
DC output voltage (nominal) & 600 & V DC \\
DC output current (100\,kW) & 167 & A \\
Combining network & Series-parallel matrix & --- \\
Substrate material & Rogers RT/Duroid 5880 & --- \\
Diode type & GaN Schottky & --- \\
Diode junction temperature (max) & 175 & $^\circ$C \\
Thermal management & Heatsink + fan array & --- \\
Waste heat (at 100\,kW output) & 13.5 & kW \\
Array mass & 15 & kg \\
Environmental rating & IP65 & --- \\
\end{longtable}

%=============================================================================
\section{Safety System Design}
%=============================================================================

\subsection{Hazard Analysis}

\begin{table}[h]
\centering
\caption{Hazard identification and mitigation.}
\begin{tabular}{p{3cm}p{3cm}p{2cm}p{4cm}}
\toprule
Hazard & Scenario & Severity & Mitigation \\
\midrule
RF exposure (personnel) & Person enters corridor & High & Keep-out volume, radar detection, auto-shutdown \\
RF exposure (wildlife) & Bird enters corridor & Medium & Acoustic deterrents, corridor height control \\
Equipment failure (TX) & Gyrotron arc & Medium & Crowbar circuit, fault detection, auto-shutdown \\
Equipment failure (RX) & Rectenna element failure & Low & Graceful degradation, monitoring \\
Corridor collapse & $\psi$-field loss & Low & Natural dispersion (safe), power ramp-down \\
EMI to nearby systems & Sidelobe radiation & Low & Corridor confinement, $<$-40\,dB sidelobes \\
Cyber attack & Unauthorized power diversion & Medium & Encrypted C3, physical interlocks \\
\bottomrule
\end{tabular}
\end{table}

\subsection{Keep-Out Volume Specification}

For 100\,kW at 94\,GHz through a corridor with $w_0 = 0.13\,\text{m}$:

\paragraph{On-axis power density:}
\begin{equation}
S_0 = \frac{P}{\pi w_0^2} = \frac{10^5}{\pi \times 0.0169}
= 1.88 \times 10^6\,\text{W/m}^2
\end{equation}

\paragraph{MPE limits (IEEE C95.1, general public, 94\,GHz):}
$S_{\rm MPE} = 10\,\text{W/m}^2$.

\paragraph{Keep-out radius:}
\begin{equation}
\rho_{\rm KO} = w_0\sqrt{\frac{1}{2}\ln\frac{S_0}{S_{\rm MPE}}}
= 0.13\sqrt{\frac{1}{2}\ln(1.88 \times 10^5)}
= 0.13 \times 2.44 = 0.32\,\text{m}
\end{equation}

\paragraph{Safety margin:} Apply 3$\times$ margin:
$\rho_{\rm KO,design} = 1.0\,\text{m}$ from corridor center.

\paragraph{Keep-out volume:} Cylinder of radius 1.0\,m, length 1\,km,
centered on corridor axis. Total volume: $\pi \times 1^2 \times 1000
= 3{,}142\,\text{m}^3$.

\subsection{Intrusion Detection System}

\begin{itemize}
\item \textbf{Primary sensor:} 94\,GHz FMCW radar (shared frequency band
with power beam), range resolution 0.15\,m, update rate 100\,Hz.
Detection range: corridor length + 100\,m buffer.

\item \textbf{Secondary sensor:} 1550\,nm LIDAR array (eye-safe), 360$^\circ$
coverage every 50\,m along corridor. Object classification capability
(distinguish birds from aircraft from personnel).

\item \textbf{Response time:} Intrusion detection to beam shutoff: $<$1\,ms.

\item \textbf{Beam shutoff mechanism:} Solid-state RF switch (PIN diode) in
waveguide feed, insertion loss 0.3\,dB, switching time 50\,ns.
Backup: crowbar circuit on gyrotron power supply, switching time 10\,$\mu$s.
\end{itemize}

\subsection{Fail-Safe Modes}

\begin{enumerate}
\item \textbf{Corridor obstruction detected:} Immediate power shutoff
($<$1\,ms). Power held at zero until obstruction clears and 5\,s
clear period confirmed.

\item \textbf{RX communication loss:} Power reduced to 10\% within 100\,ms.
Full shutoff if communication not restored within 10\,s.

\item \textbf{$\psi$-corridor degradation:} Corridor mode quality continuously
monitored via pilot tone. If mode quality drops below 80\% threshold,
power reduced proportionally. Below 50\%, full shutoff.

\item \textbf{Grid fault (TX side):} UPS maintains corridor former power
for 30\,s (allowing graceful load transfer). RF power ramps down
immediately.

\item \textbf{Grid fault (RX side):} RX signals TX to reduce power.
Excess power absorbed by RX dummy load (rated for 100\% power,
30\,s thermal capacity).
\end{enumerate}

%=============================================================================
\section{Installation and Commissioning}
%=============================================================================

\subsection{Site Requirements}

\paragraph{TX station:}
\begin{itemize}
\item Floor area: 10\,m $\times$ 20\,m (equipment room + corridor former)
\item Ceiling height: 5\,m (for gyrotron and waveguide routing)
\item Grid connection: 750\,kVA, 3-phase, 480\,V
\item Cooling: 500\,kW cooling capacity (water or air)
\item Foundation: Standard industrial, 500\,kg/m$^2$ floor loading
\end{itemize}

\paragraph{RX station:}
\begin{itemize}
\item Floor area: 5\,m $\times$ 5\,m (rectenna + power conditioning)
\item Grid connection: 200\,kVA (output)
\item Cooling: 50\,kW (waste heat from rectification)
\end{itemize}

\paragraph{Corridor path:}
\begin{itemize}
\item Clear line of sight between TX and RX
\item Keep-out volume maintained (1\,m radius cylinder)
\item No permanent structures within 3\,m of corridor center
\item Corridor height: minimum 5\,m above ground (adjustable via steering)
\end{itemize}

\subsection{Commissioning Procedure}

\begin{enumerate}
\item \textbf{Phase 1: Equipment installation and checkout} (2 weeks)
    \begin{itemize}
    \item Install TX and RX stations
    \item Connect power, cooling, communications
    \item Individual subsystem verification
    \end{itemize}

\item \textbf{Phase 2: Corridor formation} (1 week)
    \begin{itemize}
    \item Cool SRF cavities to operating temperature
    \item Establish TE+TM modes in all cavities
    \item Phase-lock cavity array
    \item Begin $2\omega$ modulation; monitor for $\psi$-mode buildup
    \item Verify corridor profile with probe beam
    \end{itemize}

\item \textbf{Phase 3: Low-power commissioning} (1 week)
    \begin{itemize}
    \item Inject 1\,kW probe signal through corridor
    \item Measure received power and mode profile
    \item Verify coupling efficiency and corridor loss
    \item Test safety interlocks with simulated intrusions
    \end{itemize}

\item \textbf{Phase 4: Power ramp-up} (1 week)
    \begin{itemize}
    \item Increase power in 10\,kW steps
    \item Monitor all thermal, electrical, and safety parameters
    \item Verify corridor stability under load
    \item Final acceptance testing at rated power
    \end{itemize}
\end{enumerate}

Total commissioning time: 5 weeks from equipment delivery.

%=============================================================================
\section{Estimated Costs and Schedule}
%=============================================================================

\subsection{Capital Cost Estimate}

\begin{table}[h]
\centering
\caption{Capital cost estimate for 1\,km, 100\,kW reference system.}
\begin{tabular}{lrr}
\toprule
Item & Quantity & Cost (USD) \\
\midrule
94\,GHz Gyrotron + power supply & 1 & \$500{,}000 \\
SRF cavity array (100 cavities) & 1 & \$2{,}000{,}000 \\
Cryoplant (2\,kW at 2\,K) & 1 & \$500{,}000 \\
Cryostat + distribution & 1 & \$300{,}000 \\
TX horn antenna & 1 & \$50{,}000 \\
Rectenna array + housing & 1 & \$200{,}000 \\
Power conditioning unit & 1 & \$100{,}000 \\
Safety system (radar + LIDAR + interlocks) & 1 & \$300{,}000 \\
C3 system & 1 & \$100{,}000 \\
Installation + commissioning & 1 & \$500{,}000 \\
\midrule
\textbf{Total capital cost} & & \textbf{\$4{,}550{,}000} \\
\bottomrule
\end{tabular}
\end{table}

\subsection{Operating Cost Estimate}

\begin{table}[h]
\centering
\caption{Annual operating cost estimate.}
\begin{tabular}{lr}
\toprule
Item & Annual Cost (USD) \\
\midrule
Electricity (581\,kW $\times$ 8760\,h $\times$ \$0.10/kWh) & \$508{,}000 \\
Helium replenishment (5\% annual loss) & \$25{,}000 \\
Gyrotron tube replacement (every 2 years) & \$100{,}000 \\
Maintenance labor (2 FTE) & \$200{,}000 \\
Insurance and permits & \$50{,}000 \\
\midrule
\textbf{Total annual operating cost} & \textbf{\$883{,}000} \\
\bottomrule
\end{tabular}
\end{table}

\subsection{Levelized Cost of Energy Transfer}

For 100\,kW delivered, 95\% availability (8322 hours/year):
\begin{equation}
\text{LCOE}_{\rm transfer} = \frac{\text{Annualized capital} + \text{O\&M}}
{\text{Energy delivered}} = \frac{455{,}000 + 883{,}000}{832{,}200\,\text{kWh}}
= \$1.61/\text{kWh}
\end{equation}
(assuming 10-year capital recovery).

This is significantly higher than conventional power lines (\$0.01--0.05/kWh
for overhead lines). The $\psi$-corridor system is economically viable
primarily for:
\begin{itemize}
\item Locations where physical cables are impossible (across water,
environmentally protected areas)
\item Emergency/rapid deployment (disaster relief)
\item Mobile power delivery (military forward bases)
\item Space applications (no alternative exists)
\end{itemize}

At scale (10\,MW+ systems), the LCOE drops to $\sim$\$0.20/kWh due to
the better utilization of the fixed-cost corridor former.

%=============================================================================
\section{Technology Roadmap}
%=============================================================================

\begin{table}[h]
\centering
\caption{Technology development roadmap.}
\begin{tabular}{llll}
\toprule
Phase & Timeline & Milestone & TRL \\
\midrule
0 & Year 0--1 & Measure $|\lambda - 1|$ in laboratory & 1$\to$2 \\
1 & Year 1--2 & Demonstrate $\psi$-mode excitation & 2$\to$3 \\
2 & Year 2--4 & Lab-scale corridor (1\,m) with guided mode & 3$\to$4 \\
3 & Year 4--6 & 100\,m outdoor corridor, 10\,kW & 4$\to$5 \\
4 & Year 6--8 & 1\,km reference system, 100\,kW & 5$\to$6 \\
5 & Year 8--10 & Multi-km system, MW-class & 6$\to$7 \\
6 & Year 10--12 & Commercial deployment & 7$\to$9 \\
\bottomrule
\end{tabular}
\end{table}

\begin{tcolorbox}[colback=red!5!white, colframe=red!60!black, title=Critical Path Item]
\textbf{The entire roadmap is gated by Phase 0: measurement of $|\lambda - 1|$.}
If $|\lambda - 1| = 0$ exactly (no EM-$\psi$ back-reaction), then direct
$\psi$-corridor generation via EM pumping is impossible. Alternative
approaches (mass-sourced corridors, metamaterial-only systems) would then
be pursued, with significantly different timelines and capabilities.
\end{tcolorbox}

%=============================================================================
\section{Intellectual Property Coverage}
%=============================================================================

Patent \#4 covers the following claims:

\begin{enumerate}
\item A method for wireless power transfer comprising: generating a shaped
$\psi$-field envelope defining a corridor of modified refractive index between
a transmitter and a receiver; coupling electromagnetic energy into guided modes
of said corridor; propagating said guided modes along said corridor; and
extracting electromagnetic energy from said guided modes at the receiver.

\item The method of claim 1 wherein the $\psi$-field envelope is generated by
an array of electromagnetic resonant cavities operating in a superposition of
transverse electric and transverse magnetic modes with amplitude modulation at
twice the drive frequency.

\item A system for $\psi$-field wireless power transfer comprising: a transmitter
station with an RF source and $\psi$-corridor former; a $\psi$-corridor
transmission medium; a receiver station with a $\psi$-responsive coupler and
rectification system; and a safety system with keep-out volume enforcement.

\item The system of claim 3 further comprising metamaterial beamforming elements
for enhanced mode coupling.

\item The system of claim 3 further comprising multi-lane delivery using
mode-division multiplexing within a single $\psi$-corridor.

\item A method for dynamic electric vehicle charging comprising a
ground-embedded $\psi$-corridor former and vehicle-mounted receiver coupler.

\item A method for deep-space energy beaming using $\psi$-corridors established
and maintained in vacuum between orbital or planetary stations.

\item A safety system for $\psi$-corridor power transfer comprising corridor
volume surveillance, beam interlock, and automatic shutdown upon intrusion
detection.
\end{enumerate}

%=============================================================================
\section{Conclusions and Recommendations}
%=============================================================================

This engineering design document establishes that a $\psi$-corridor wireless
power transfer system is technically feasible using existing component
technologies (gyrotrons, SRF cavities, rectennas, radar safety systems),
contingent upon the existence of a nonzero EM-$\psi$ coupling parameter
$|\lambda - 1|$.

\paragraph{Key recommendations:}
\begin{enumerate}
\item \textbf{Priority 1:} Execute the $|\lambda - 1|$ measurement protocol
described in Appendix R of the main DFD text. This is the single most
important experiment for the entire $\psi$-corridor technology programme.

\item \textbf{Priority 2:} Develop the metamaterial-only $\psi$-corridor
analogue as a parallel track. This provides immediate practical value
(GRIN waveguide in air) while the $\psi$-field physics is being established.

\item \textbf{Priority 3:} Design and build a proof-of-concept corridor
former using existing SRF infrastructure (e.g., at DESY, FNAL, or
Jefferson Lab). The cavity technology is mature; only the TE+TM
configuration and $2\omega$ modulation scheme are novel.

\item \textbf{Priority 4:} Engage with the space power community
(NASA, JAXA, ESA) to evaluate $\psi$-corridor potential for space-based
solar power and lunar/Mars power delivery. The space application has
the largest potential impact and the least competition from conventional
alternatives.
\end{enumerate}

\end{document}
